% TOP_PROBLEM: {"id": 30002380, "title": "Serre's Uniformity Problem for Elliptic Curves", "category_id": 1, "category": {"id": 1, "name": "number_theory", "display_name": "Number Theory", "description": "Properties of integers, prime numbers, Diophantine equations.", "slug": "number-theory", "order_index": 1, "created_at": "2026-07-31T15:26:25.670Z"}, "status": "partially_solved"}
% FIRST_POSTED: 2026-09-25
% ENTRY_KIND: statement_only
% !TeX program = lualatex
% Definition and sources only; this file is not an LLM solution attempt.
\documentclass[11pt]{article}
\usepackage[margin=25mm]{geometry}
\usepackage{fontspec}
\setmainfont{Latin Modern Roman}
\usepackage{amsmath,amssymb,mathtools}
\usepackage{unicode-math}
\setmathfont{Latin Modern Math}
\usepackage{xurl,hyperref}
\hypersetup{colorlinks=true,urlcolor=blue}
\setcounter{secnumdepth}{0}
\setlength{\parindent}{0pt}
\setlength{\parskip}{5pt}
\setlength{\emergencystretch}{3em}
\begin{document}
\section*{\texorpdfstring{Serre's Uniformity Problem for Elliptic Curves}{Serre's Uniformity Problem for Elliptic Curves}}\label{30002380}
\subsection{Definitions and mathematical statement}
For an elliptic curve $E/{\mathbb{Q}}$ and a prime $p$, its geometric $p$-torsion is a two-dimensional ${\mathbb{F}}_p$-vector space. Choosing a basis gives
$\rho_{E,p}:\operatorname{Gal}(\overline{{\mathbb{Q}}}/{\mathbb{Q}})\to\mathrm{GL}_2({\mathbb{F}}_p)$, well defined up to conjugation. Non-CM means $\operatorname{End}_{\overline{{\mathbb{Q}}}}(E)={\mathbb{Z}}$. The question is
\[
 \exists N>0\quad\forall E/{\mathbb{Q}}\text{ non-CM}\quad\forall p>N,
 \qquad\operatorname{im}\rho_{E,p}=\mathrm{GL}_2({\mathbb{F}}_p).
\]
The bound must be independent of the curve, not merely finite for each fixed curve. It is surjectivity of the full residual representation, not just irreducibility or a condition on its determinant.
\par\smallskip\textit{Formulation and concept references:} \href{https://arxiv.org/abs/2305.17780}{[S1]}.

\subsection{Short English statement}
Is there a single prime cutoff beyond which every non-CM elliptic curve over the rationals has the largest possible Galois action on its torsion points?
\subsection{Sources}
\begin{itemize}
\item[S1]Formal statement, abstract.
\url{https://arxiv.org/abs/2305.17780}.
\textit{Formulation source.}
\end{itemize}
\end{document}
